% Canonical Paper IV appendix.
This appendix expands the residual construction, especially the partial-domain
and encoding conventions suppressed in the main flow.

\subsection{Why an invalid outcome is necessary}

Suppose a continuation class contains the ordinary future $z\mapsto1/z$ but
one compares only outputs on futures legal for both pasts.  The current value
$0$ has no legal output for this future.  A common-domain convention can then
make its comparison vacuous.  Totalization instead records
\[
  O(0\cdot(1/z))=\bot,
  \qquad
  O(1\cdot(1/z))=1,
\]
so the two pasts are correctly distinguished.  The same device handles type
errors, missing registers, and forbidden resource actions.

Associativity remains valid after totalization: once a computation enters
$\bot$, every further continuation stays there.  Therefore the proof of the
right-congruence proposition requires no special cases.

\subsection{The factorization diagram}

For an exact deterministic machine, Theorem~\ref{thm:minimal-exact-residual}
produces the commutative diagram
\[
\begin{tikzpicture}[baseline=(current bounding box.center),node distance=18mm and 25mm]
\node (P) {$P_C$};
\node[right=of P] (S) {$e(P_C)$};
\node[below=of S] (R) {$\cR_C$};
\draw[-{Latex}] (P) -- node[above] {$e$} (S);
\draw[-{Latex}] (P) -- node[below left] {$p\mapsto[p]$} (R);
\draw[-{Latex}] (S) -- node[right] {$\phi$} (R);
\end{tikzpicture}
\]
Every transition commutes with $\phi$, and the output factors through
$\cR_C$.  Thus a larger exact state representation may contain caches or
redundancy, but it cannot merge two residual classes.

\subsection{One-way communication interpretation}

At a cut, let Alice know the past $p$ and Bob choose the future $\eta$.  Alice
sends one deterministic message $m(p)$; Bob must output
$O_C(p\cdot\eta)$ for every legal pair.

\begin{proposition}[Deterministic one-way lower bound]
\label{prop:one-way-residual-bound}
Any such exact protocol uses distinct messages for distinct residual classes.
A fixed-width worst-case message therefore has length at least
$\lceil\log_2|\cR_C|\rceil$.
\end{proposition}

\begin{proof}
If $m(p)=m(p')$, Bob behaves identically for every chosen future.  Exactness
then gives $p\equiv_Cp'$.  The counting bound follows.
\end{proof}

This interpretation makes clear why the future class matters: it is Bob's set
of allowed tests.

\subsection{Dense tensor boundary example}

Let the state exposed at a cut be a dense tensor represented as a vector
$T\in\F_q^D$.  If every linear functional
$\ell\in(\F_q^D)^\vee$ is an allowed future observation, then two states are
residually equivalent exactly when they are equal.  Indeed, if $T-T'\ne0$, a
coordinate functional detects a nonzero component.  Hence
\[
  |\cR_C|=q^D,
  \qquad
  \log_2|\cR_C|=D\log_2q.
\]
A direct fixed-width residue encoding uses $D\lceil\log_2q\rceil$ bits.

If only one nonzero functional $\ell$ is allowed, the residual quotient is
instead $\F_q^D/\ker\ell$, which has $q$ classes.  The boundary dimension did
not change; the legal future class did.  Sparse, low-rank, symmetric, or
approximate tensor models change the representation again.

\subsection{Prefix-free variable-length states}

Let $N$ residual classes receive a prefix-free binary code with maximum length
$b$.  Kraft's inequality gives
\[
  1\ge\sum_{i=1}^{N}2^{-\ell_i}
  \ge N2^{-b},
\]
so $b\ge\lceil\log_2N\rceil$.  Without prefix-freeness or another decoding
convention, a raw collection of variable-length strings cannot be counted as a
self-delimiting state code.

Average code length obeys distribution-sensitive entropy bounds, while the
main theorem is a worst-case state-selection statement.  The two should not be
interchanged.

\subsection{Infinite residual spaces}

When $\cR_C$ is infinite, cardinal logarithms do not directly give finite bit
bounds.  One may need:
\begin{itemize}
  \item a topology and covering or packing numbers at scale $\varepsilon$;
  \item a measure and Shannon or mutual information;
  \item a metric loss and rate--distortion function;
  \item an algebraic dimension or rank;
  \item a computable presentation and description-length convention.
\end{itemize}
No one of these is canonical without an application-specific observation and
approximation criterion.

\subsection{Changing future classes}

Let $M_1\subseteq M_2$ be future classes with the same pasts and observation.
Then
\[
  p\equiv_{M_2}p'
  \Longrightarrow
  p\equiv_{M_1}p'.
\]
Thus there is a natural surjection
\[
  P/{\equiv_{M_2}}\longrightarrow P/{\equiv_{M_1}}.
\]
Adding tests refines the residual space; removing tests condenses it.  The
projective passage from left-only futures to arbitrary two-sided futures is a
group-theoretic instance of this general refinement.
